% Generated draft: requires mathematical and source review.
\documentclass[11pt]{article}
\usepackage[margin=25mm]{geometry}
\usepackage{fontspec}
\setmainfont{Latin Modern Roman}
\usepackage{amsmath,amssymb,mathtools,unicode-math}
\setmathfont{Latin Modern Math}
\usepackage{enumitem,needspace,xurl,hyperref}
\setlength{\parindent}{0pt}
\setlength{\parskip}{5pt}
\setlength{\emergencystretch}{3em}
\newcommand{\N}{\mathbb{N}}
\newcommand{\Z}{\mathbb{Z}}
\newcommand{\Q}{\mathbb{Q}}
\newcommand{\R}{\mathbb{R}}
\newcommand{\C}{\mathbb{C}}
\newcommand{\F}{\mathbb{F}}
\newcommand{\A}{\mathbb{A}}
\newcommand{\PP}{\mathbb{P}}
\newcommand{\E}{\mathbb{E}}
\newcommand{\eps}{\varepsilon}
\newcommand{\dd}{\,\mathrm{d}}
\begin{document}
\section*{UnsolvedMath ID: 30004813}
\textbf{Source problem code:} catalogue entry
\subsection{Definitions and mathematical statement}
For a closed, connected, oriented three-manifold \(M\), irreducibility means that every embedded two-sphere bounds a three-ball. For a tame knot \(K\subset M\), let \(\nu K\) be a closed tubular neighbourhood and define its exterior by
\[
E(M,K)=M\setminus\operatorname{int}(\nu K),
\]
with the orientation induced from \(M\). Its open complement is \(M\setminus K\). The tameness assumption makes the exterior notation meaningful; the supplied record does not specify the category of knots or whether ``complement'' means exterior or open complement. No orientation on the knot itself is specified.

One possible meaning of equivalence is orientation-preserving equivalence of pairs:
\[
(M,K)\cong_{+}(N,L)
\quad\Longleftrightarrow\quad
\text{there exists an orientation-preserving homeomorphism }F:M\longrightarrow N
\text{ with }F(K)=L.
\]
This need not mean ambient isotopy, and it does not require a specified homeomorphism of exteriors to extend.

\textbf{Complement question, with unresolved scope.}
Determine the appropriate hypotheses and notion of knot equivalence for a Gordon--Luecke analogue in general three-manifolds, and determine whether that analogue holds. A precise candidate suggested by the example in the record is the following implication for every closed, connected, oriented, irreducible three-manifold \(M\) and every pair of tame knots \(K,L\subset M\) satisfying the intended exclusions:
\[
E(M,K)\cong_{+}E(M,L)
\quad\Longrightarrow\quad
(M,K)\cong_{+}(M,L).
\]
Here the premise means existence of an orientation-preserving homeomorphism of exteriors, with no boundary marking imposed. This is only a candidate interpretation: the record does not settle whether the ambient manifolds must coincide, whether equivalence means equivalence of pairs or ambient isotopy, or whether a given complement homeomorphism must extend. The phrase ``outside the rational-unknot lens-space exceptions'' does not specify the excluded knot pairs. Consequently it does not determine the domain of the universal quantifier above. The record also does not define its ``strong'' and ``weak'' analogues, so those labels cannot be assigned to this displayed candidate from the supplied evidence alone.

\textbf{Prime-decomposition question, with unresolved conventions.}
Determine what prime-decomposition theory holds for knots in general three-manifolds. Under the usual knot-pair connected-sum convention, \((M_1,K_1)\mathbin{\#}(M_2,K_2)\) is formed by removing a three-ball meeting each knot in a trivial arc and gluing the resulting sphere boundaries compatibly with orientations and arc endpoints. The identity pair is \((S^3,U)\), where \(U\) is the unknot. A nonidentity pair is prime under this convention if every expression as a connected sum has an identity factor. Possible assertions to distinguish are existence of a finite expression
\[
(M,K)\cong_{+}(M_1,K_1)\mathbin{\#}\cdots\mathbin{\#}(M_r,K_r),
\qquad r\in\N,
\]
with prime factors, uniqueness up to permutation and the intended equivalence, and termination of successive decompositions. These are separate properties. The source does not specify the ambient class or decomposition conventions sufficiently to select a single precise assertion.

The supplied dated curation labels the combined record partially solved, reports that a strong analogue is false and a weak analogue remains open in the stated generality, and attributes an answer to the separate prime-decomposition question to Miyazaki, with nonterminating decompositions possible. These are supplied literature assessments, not independently verified mathematical conclusions here; in particular, nontermination alone does not specify which existence or uniqueness assertion is intended.
\subsection{Short English statement}
The problem asks when an orientation-preserving homeomorphism of knot complements determines the knot, and what existence, uniqueness, or termination properties hold for prime decompositions of knots in general three-manifolds. The exceptional knots, equivalence relation, and strong/weak formulations are insufficiently specified. Supplied dated curation reports partial resolution; this draft requires human mathematical review.
\subsection{Sources}
\begin{itemize}
\item Local UnsolvedMath registry, permanent ID \(30004813\), exact source code \texttt{catalogue entry}, title: \emph{Knot Complements and Prime Decomposition in General Three-Manifolds}; collection: Oberwolfach Reports Open Problems. No individual proposer is supplied; the proposed year is recorded as 2021.
\item Supplied source citation: \emph{MATRIX-MFO Tandem Workshop: Invariants and Structures in Low-Dimensional Topology (2022)}, Oberwolfach Reports, pp.\ 2327--2364, DOI 10.4171/owr/2021/42. Source URL and literature-source link labelled ``OWR report'': \url{https://doi.org/10.4171/owr/2021/42}. The report title gives 2022, while the DOI and registry year refer to 2021; this discrepancy remains unresolved.
\item Supplied literature-source link labelled ``Aceto et al.'': \url{https://doi.org/10.1112/jlms.70600}. The record supplies no article title, complete author list, or theorem locator.
\item Literature assessment dated exactly \texttt{2026-08-22}: supplied dated curation, not independently verified. It reports that the strong Gordon--Luecke analogue is false, the weak analogue remains open in the stated generality, and the separate prime-decomposition question was answered by Miyazaki, with nonterminating decompositions possible. No specific Miyazaki citation is supplied. The background explicitly requests independent verification of time-sensitive status and cited claims.
\item The registry labels the statement \texttt{corrected\_verified} and describes its verification as a publication-format audit preserving mathematical content while repairing notation or typography. This supplied label does not establish mathematical correctness or resolve the formulation ambiguities.
\end{itemize}
\end{document}
